\documentclass[10pt,a4paper]{amsart}
\usepackage[margin=19mm]{geometry}
\usepackage{amsmath,amssymb,mathtools}
\usepackage[scr=rsfs,cal=euler]{mathalfa}
\usepackage{xfrac}
\usepackage[unicode,hidelinks,bookmarks=false]{hyperref}
\newcommand{\ie}{i.e.\ }
\newcommand{\cf}{cf.\ }
\newcommand{\resp}{respectively\ }
\newcommand{\Subsection}{\subsection}
\providecommand{\nobreakdash}{\nobreak\hbox{-}}
\setlength{\emergencystretch}{2em}
\multlinegap0pt
\usepackage{amssymb}
\usepackage{csquotes}
\usepackage{tikz}
\usepackage{enumerate}
\newtheorem{prop}{Proposition}
\title{SU(n)-structures through quotient by torus actions}
\author{Quentin Peres}
\date{Source: arXiv:2511.05257v2 (CC BY 4.0)}
\begin{document}
\maketitle

\noindent\emph{Source and licence.} SU(n)-structures through quotient by torus actions. Version arXiv:2511.05257v2 (CC BY 4.0), distributed by arXiv under the Creative Commons Attribution 4.0 International licence, http://creativecommons.org/licenses/by/4.0/, as recorded in the arXiv licence metadata for this exact version and shown on the abstract page https://arxiv.org/abs/2511.05257v2. Full licence notice: \texttt{licenses/arxiv-2511.05257-CC-BY-4.0.txt}.\par\smallskip

\noindent\textit{Editorial scope.} Counted unit: the article's prop eqintermediaire1 (source0.tex lines 182--187). The article's complete original proof of it is delivered with it (source0.tex lines 188--190, one \texttt{proof} environment of 3 lines). No mathematics is written, repaired or re-derived: the statement and the proof are the article's own lines, in the article's own notation, and no retained line is substituted or re-flowed.  No further source-local context is needed: every cross-reference of every retained line resolves inside the two delivered spans.  Nothing was omitted from any retained span, so the package carries no omission record.  No retained line contains a citation, so the wrapper carries no bibliography and no bibliography entry.  No retained line contains a figure environment, an image-inclusion command or drawing code, so no image is delivered and the package declares no asset dependency.  Editorial formatting only, in addition: an AMS \texttt{amsart} preamble that restates the article's own declarations and macros used by the retained lines, namely the environments \texttt{prop} on separate counters, together with the standard packages those lines need.  The retained environments print in this document as Proposition 1 in order of appearance, whereas the article prints the same items with the numbering of its own full document.  The complete native source of the article as distributed in the arXiv e-print for this exact version is bundled unchanged as \texttt{sources/188648/source0.tex}.
\begin{prop} $\Omega_M$ and $\omega_M$ satisfy the following equations:
    \begin{align}
        \Omega_M\wedge\omega_M&=0, \label{eqintermediaire1} \\
        \Omega_M\wedge\overline{\Omega_M}&=c_n\|V_1\|^2\cdots\|V_s\|^2\omega_M^n.\label{eqintermediaire2}
    \end{align}
\end{prop}

\begin{proof}
    Using $\Omega_X\wedge\omega_X=0$ and contracting successively by $\xi_1,\cdots\xi_s$, we obtain equation (\ref{eqintermediaire1}). For equation  (\ref{eqintermediaire2}), start from $\Omega_X\wedge\overline{\Omega_X}=c_{n+s}\omega_X^{n+s}$ and contract with $\xi_1,\overline{\xi_1},\cdots,\xi_s,\overline{\xi_s}$. Noticing that $\iota_{\overline{\xi_s}}\Omega_X=0$ and $\omega_X(\xi_k,\overline{\xi_k})=\frac{i}{2}\|V_k\|^2$, we obtain the second equation.
\end{proof}

\end{document}
